% GENERATED FILE. Edit canonical lessons, then run npm run generate:books.
% canonical-source-hash: fnv1a64:28d9617c1f274739
% canonical-lessons: SA-C502-karmasala, SA-C502-taksakah, SA-C502-taksati, SA-C502-krakacah, SA-C502-tankah

\chapter{In the Workshop}
\label{ch:sa-in-the-workshop}
\hlchaptermodality{502}

\begin{chapteropening}
\textbf{By the end of this chapter:} \emph{I can say where a craftsman works and which tools he uses.}

Complete the last lesson of chapter 502: I can say where a craftsman works and which tools he uses.
\end{chapteropening}

\section[\texorpdfstring{\sk{कर्मशाला} (karmaśālā)}{कर्मशाला (karmaśālā)}]{\sk{कर्मशाला} (karmaśālā) — a workshop}

\label{lesson:SA-C502-karmasala}



\begin{quote}
Before the new one: say the Sanskrit for radio, then the Sanskrit for the internet.
\end{quote}

\subsection*{You'll want to know: \sk{कर्मशाला}}
\textbf{\sk{कर्मशाला}} — \emph{karmaśālā} — \textquotedblleft{}a workshop\textquotedblright{}.

\begin{quote}
\sk{मम} \sk{पिता} \sk{कर्मशालायां} \sk{कार्यं} \sk{करोति।}
\end{quote}

\begin{quote}
\emph{mama pitā karmaśālāyāṃ kāryaṃ karoti.}
\end{quote}

\begin{quote}
My father works in a workshop.
\end{quote}

A feminine noun.

\begin{cousinweb}
\textbf{\sk{कर्मशाला}} (\emph{karmaśālā}) is built from \emph{karman}, \textquotedblleft{}work\textquotedblright{}, and \emph{śālā}, \textquotedblleft{}a hall, a room\textquotedblright{}; MW records \emph{karmaśālā} as the room where the day's work is done, a workshop (MW).
\end{cousinweb}

\begin{grammarlens}[title={What you've built}]
One more word for describing what you lived and felt.
\end{grammarlens}

\subsection*{Your turn}
\begin{itemize}
\raggedright
  \item \emph{Say it:} \emph{karmaśālā}
  \item \emph{Say it:} \emph{karmaśālā}, once more
  \item \emph{Say it:} \emph{mama pitā karmaśālāyāṃ kāryaṃ karoti.} — My father works in a workshop.
  \item \emph{Recall:} say the Sanskrit for radio, then the Sanskrit for the internet, then say \emph{karmaśālā} again
\end{itemize}

\begin{tcolorbox}[breakable,colback=teal!4,colframe=teal!35!black,title={Before you move on}]
What does \sk{कर्मशाला} mean? (\textquotedblleft{}A workshop\textquotedblright{}.) Say it once more.
\end{tcolorbox}
\section[\texorpdfstring{\sk{तक्षकः} (takṣakaḥ)}{तक्षकः (takṣakaḥ)}]{\sk{तक्षकः} (takṣakaḥ) — a carpenter}

\label{lesson:SA-C502-taksakah}



\begin{quote}
Before the new one: say the Sanskrit for the internet, then the Sanskrit for a workshop.
\end{quote}

\subsection*{You'll want to know: \sk{तक्षकः}}
\textbf{\sk{तक्षकः}} — \emph{takṣakaḥ} — \textquotedblleft{}a carpenter\textquotedblright{}.

\begin{quote}
\sk{तक्षकः} \sk{नूतनाम्} \sk{आसन्दीं} \sk{रचयति।}
\end{quote}

\begin{quote}
\emph{takṣakaḥ nūtanām āsandīṃ racayati.}
\end{quote}

\begin{quote}
The carpenter makes a new chair.
\end{quote}

A masculine noun.

\begin{cousinweb}
\textbf{\sk{तक्षकः}} (\emph{takṣakaḥ}) comes from the root √takṣ, \textquotedblleft{}to pare, to fashion by cutting\textquotedblright{}, with the agent ending \emph{-aka}; MW gives \emph{takṣaka} the sense \textquotedblleft{}a carpenter\textquotedblright{} (MW).
\end{cousinweb}

\begin{grammarlens}[title={What you've built}]
One more word for describing what you lived and felt.
\end{grammarlens}

\subsection*{Your turn}
\begin{itemize}
\raggedright
  \item \emph{Say it:} \emph{takṣakaḥ}
  \item \emph{Say it:} \emph{takṣakaḥ}, once more
  \item \emph{Say it:} \emph{takṣakaḥ nūtanām āsandīṃ racayati.} — The carpenter makes a new chair.
  \item \emph{Recall:} say the Sanskrit for the internet, then the Sanskrit for a workshop, then say \emph{takṣakaḥ} again
\end{itemize}

\begin{tcolorbox}[breakable,colback=teal!4,colframe=teal!35!black,title={Before you move on}]
What does \sk{तक्षकः} mean? (\textquotedblleft{}A carpenter\textquotedblright{}.) Say it once more.
\end{tcolorbox}
\section[\texorpdfstring{\sk{तक्षति} (takṣati)}{तक्षति (takṣati)}]{\sk{तक्षति} (takṣati) — pares; shapes wood}

\label{lesson:SA-C502-taksati}



\begin{quote}
Before the new one: say the Sanskrit for a workshop, then the Sanskrit for a carpenter.
\end{quote}

\subsection*{You'll want to know: \sk{तक्षति}}
\textbf{\sk{तक्षति}} — \emph{takṣati} — \textquotedblleft{}pares; shapes wood\textquotedblright{}.

\begin{quote}
\sk{तक्षकः} \sk{दारु} \sk{तक्षति।}
\end{quote}

\begin{quote}
\emph{takṣakaḥ dāru takṣati.}
\end{quote}

\begin{quote}
The carpenter shapes the wood.
\end{quote}

Past (he …): \sk{सः} \sk{दारु} \sk{तष्टवान्}, \emph{saḥ dāru taṣṭavān}.

\begin{cousinweb}
\textbf{\sk{तक्षति}} (\emph{takṣati}) comes from the root √takṣ, \textquotedblleft{}to form by cutting, to chisel\textquotedblright{}; its past participle is \emph{taṣṭa}, and KRM lists the past \emph{taṣṭavān} used here (MW).
\end{cousinweb}

\begin{grammarlens}[title={What you've built}]
One more word for describing what you lived and felt.
\end{grammarlens}

\subsection*{Your turn}
\begin{itemize}
\raggedright
  \item \emph{Say it:} \emph{takṣati}
  \item \emph{Say it:} \emph{takṣati}, once more
  \item \emph{Say it:} \emph{takṣakaḥ dāru takṣati.} — The carpenter shapes the wood.
  \item \emph{Recall:} say the Sanskrit for a workshop, then the Sanskrit for a carpenter, then say \emph{takṣati} again
\end{itemize}

\begin{tcolorbox}[breakable,colback=teal!4,colframe=teal!35!black,title={Before you move on}]
What does \sk{तक्षति} mean? (\textquotedblleft{}Pares; shapes wood\textquotedblright{}.) Say it once more.
\end{tcolorbox}
\section[\texorpdfstring{\sk{क्रकचः} (krakacaḥ)}{क्रकचः (krakacaḥ)}]{\sk{क्रकचः} (krakacaḥ) — a saw}

\label{lesson:SA-C502-krakacah}



\begin{quote}
Before the new one: say the Sanskrit for a carpenter, then the Sanskrit for pares.
\end{quote}

\subsection*{You'll want to know: \sk{क्रकचः}}
\textbf{\sk{क्रकचः}} — \emph{krakacaḥ} — \textquotedblleft{}a saw\textquotedblright{}.

\begin{quote}
\sk{सः} \sk{क्रकचेन} \sk{दारु} \sk{छिनत्ति।}
\end{quote}

\begin{quote}
\emph{saḥ krakacena dāru chinatti.}
\end{quote}

\begin{quote}
He cuts the wood with a saw.
\end{quote}

A masculine noun.

\begin{cousinweb}
\textbf{\sk{क्रकचः}} (\emph{krakacaḥ}) is the word MW records for a saw; MW gives no derivation, so its origin is uncertain (MW).
\end{cousinweb}

\begin{grammarlens}[title={What you've built}]
One more word for describing what you lived and felt.
\end{grammarlens}

\subsection*{Your turn}
\begin{itemize}
\raggedright
  \item \emph{Say it:} \emph{krakacaḥ}
  \item \emph{Say it:} \emph{krakacaḥ}, once more
  \item \emph{Say it:} \emph{saḥ krakacena dāru chinatti.} — He cuts the wood with a saw.
  \item \emph{Recall:} say the Sanskrit for a carpenter, then the Sanskrit for pares, then say \emph{krakacaḥ} again
\end{itemize}

\begin{tcolorbox}[breakable,colback=teal!4,colframe=teal!35!black,title={Before you move on}]
What does \sk{क्रकचः} mean? (\textquotedblleft{}A saw\textquotedblright{}.) Say it once more.
\end{tcolorbox}
\section[\texorpdfstring{\sk{टंकः} (ṭaṅkaḥ)}{टंकः (ṭaṅkaḥ)}]{\sk{टंकः} (ṭaṅkaḥ) — a chisel}

\label{lesson:SA-C502-tankah}



\begin{quote}
Before the new one: say the Sanskrit for pares, then the Sanskrit for a saw.
\end{quote}

\subsection*{You'll want to know: \sk{टंकः}}
\textbf{\sk{टंकः}} — \emph{ṭaṅkaḥ} — \textquotedblleft{}a chisel\textquotedblright{}.

\begin{quote}
\sk{शिल्पी} \sk{टंकेन} \sk{शिलां} \sk{तक्षति।}
\end{quote}

\begin{quote}
\emph{śilpī ṭaṅkena śilāṃ takṣati.}
\end{quote}

\begin{quote}
The sculptor shapes the stone with a chisel.
\end{quote}

A masculine noun.

\begin{cousinweb}
\textbf{\sk{टंकः}} (\emph{ṭaṅkaḥ}) is the word MW gives for a spade, a hatchet and a stone-cutter's chisel; MW gives no derivation, so its origin is uncertain (MW).
\end{cousinweb}

\begin{grammarlens}[title={What you've built}]
One more word for describing what you lived and felt.
\end{grammarlens}

\subsection*{Your turn}
\begin{itemize}
\raggedright
  \item \emph{Say it:} \emph{ṭaṅkaḥ}
  \item \emph{Say it:} \emph{ṭaṅkaḥ}, once more
  \item \emph{Say it:} \emph{śilpī ṭaṅkena śilāṃ takṣati.} — The sculptor shapes the stone with a chisel.
  \item \emph{Recall:} say the Sanskrit for pares, then the Sanskrit for a saw, then say \emph{ṭaṅkaḥ} again
\end{itemize}

\begin{tcolorbox}[breakable,colback=teal!4,colframe=teal!35!black,title={Before you move on}]
What does \sk{टंकः} mean? (\textquotedblleft{}A chisel\textquotedblright{}.) Say it once more.
\end{tcolorbox}
